\title{2017 年高考数学 全国I卷-文科}
\date{}
\maketitle

\section*{一、选择题}
本大题共12小题，每小题5分，共60分。在每小题给出的四个选项中，只有一项是符合题目要求的。

\begin{question}[points=5]
已知集合$A=\{x|x<2\}$，$B=\{x|3-2x>0\}$，则（        ）
\begin{choices}
  \item $A\cap B=\{x|x<\frac{3}{2}\}$
  \item $A\cap B=\varnothing$
  \item $A\cup B=\{x|x<\frac{3}{2}\}$
  \item $A\cup B=\mathbb{R}$
\end{choices}
\begin{solution}
\textbf{答案：}A

\textbf{分析：}解不等式求出集合$B$，结合集合交集和并集的定义，可得结论。

\textbf{详解：}解：$\because$集合$A=\{x|x<2\}$，$B=\{x|3-2x>0\}=\{x|x<\frac{3}{2}\}$，$\therefore A\cap B=\{x|x<\frac{3}{2}\}$，故A正确，B错误；$A\cup B=\{x|x<2\}$，故C，D错误；故选：A。
\end{solution}
\end{question}

\begin{question}[points=5]
为评估一种农作物的种植效果，选了$n$块地作试验田．这$n$块地的亩产量（单位：$\text{kg}$）分别是$x_1$，$x_2$，$\ldots$，$x_n$，下面给出的指标中可以用来评估这种农作物亩产量稳定程度的是（           ）
\begin{choices}
  \item $x_1$，$x_2$，$\ldots$，$x_n$的平均数
  \item $x_1$，$x_2$，$\ldots$，$x_n$的标准差
  \item $x_1$，$x_2$，$\ldots$，$x_n$的最大值
  \item $x_1$，$x_2$，$\ldots$，$x_n$的中位数
\end{choices}
\begin{solution}
\textbf{答案：}B

\textbf{分析：}利用平均数、标准差、最大值、中位数的定义和意义直接求解。

\textbf{详解：}解：在A中，平均数是表示一组数据集中趋势的量数，它是反映数据集中趋势的一项指标，故A不可以用来评估这种农作物亩产量稳定程度；在B中，标准差能反映一个数据集的离散程度，故B可以用来评估这种农作物亩产量稳定程度；在C中，最大值是一组数据最大的量，故C不可以用来评估这种农作物亩产量稳定程度；在D中，中位数将数据分成前半部分和后半部分，用来代表一组数据的“中等水平”，故D不可以用来评估这种农作物亩产量稳定程度。故选：B。
\end{solution}
\end{question}

\begin{question}[points=5]
下列各式的运算结果为纯虚数的是（                ）
\begin{choices}
  \item $\text{i}(1+\text{i})^2$
  \item $\text{i}^2(1-\text{i})$
  \item $(1+\text{i})^2$
  \item $\text{i}(1+\text{i})$
\end{choices}
\begin{solution}
\textbf{答案：}C

\textbf{分析：}利用复数的运算法则、纯虚数的定义即可判断出结论。

\textbf{详解：}解：A．$\text{i}(1+\text{i})^2=\text{i}\cdot2\text{i}=-2$，是实数．B．$\text{i}^2(1-\text{i})=-1+\text{i}$，不是纯虚数．C．$(1+\text{i})^2=2\text{i}$为纯虚数．D．$\text{i}(1+\text{i})=\text{i}-1$不是纯虚数．故选：C。
\end{solution}
\end{question}

\begin{question}[points=5]
如图，正方形$ABCD$内的图形来自中国古代的太极图．正方形内切圆中的黑色部分和白色部分关于正方形的中心成中心对称．在正方形内随机取一点，则此点取自黑色部分的概率是（                 ）

\begin{center}
  \includegraphics[width=.38\linewidth]{question-4-1.jpg}
\end{center}
\begin{choices}
  \item $\frac{1}{4}$
  \item $\frac{1}{8}$
  \item $\frac{1}{2}$
  \item $\frac{\pi}{4}$
\end{choices}
\begin{solution}
\textbf{答案：}B

\textbf{分析：}根据图象的对称性求出黑色图形的面积，结合几何概型的概率公式进行求解即可。

\textbf{详解：}解：根据图象的对称性知，黑色部分为圆面积的一半，设圆的半径为1，则正方形的边长为2，则黑色部分的面积$S=\frac{\pi}{2}$，则对应概率$P=\frac{\frac{\pi}{2}}{4}=\frac{\pi}{8}$，故选：B。
\end{solution}
\end{question}

\begin{question}[points=5]
已知$F$是双曲线$C:x^2-\frac{y^2}{3}=1$的右焦点，$P$是$C$上一点，且$PF$与$x$轴垂直，点$A$的坐标是$(1,3)$，则$\triangle APF$的面积为（             ）
\begin{choices}
  \item $\frac{1}{3}$
  \item $\frac{1}{2}$
  \item $\frac{2}{3}$
  \item $\frac{3}{2}$
\end{choices}
\begin{solution}
\textbf{答案：}D

\textbf{分析：}由题意求得双曲线的右焦点$F(2,0)$，由$PF$与$x$轴垂直，代入即可求得$P$点坐标，根据三角形的面积公式，即可求得$\triangle APF$的面积。

\textbf{详解：}解：由双曲线$C:x^2-\frac{y^2}{3}=1$的右焦点$F(2,0)$，$PF$与$x$轴垂直，设$(2,y)$，$y>0$，则$y=3$，则$P(2,3)$，$\therefore AP\perp PF$，则$|AP|=1$，$|PF|=3$，$\therefore\triangle APF$的面积$S=\frac{1}{2}\times|AP|\times|PF|=\frac{3}{2}$，同理当$y<0$时，则$\triangle APF$的面积$S=\frac{3}{2}$，故选：D。
\end{solution}
\end{question}

\begin{question}[points=5]
如图，在下列四个正方体中，$A$，$B$为正方体的两个顶点，$M$，$N$，$Q$为所在棱的中点，则在这四个正方体中，直线$AB$与平面$MNQ$不平行的是（     ）
\begin{choices}
  \item 
  \item 
  \item 
  \item 
\end{choices}
\begin{solution}
\textbf{答案：}A

\textbf{分析：}利用线面平行判定定理可知B、C、D均不满足题意，从而可得答案。

\textbf{详解：}解：对于选项B，由于$AB\parallel MQ$，结合线面平行判定定理可知B不满足题意；对于选项C，由于$AB\parallel MQ$，结合线面平行判定定理可知C不满足题意；对于选项D，由于$AB\parallel NQ$，结合线面平行判定定理可知D不满足题意；所以选项A满足题意，故选：A。
\end{solution}
\end{question}

\begin{question}[points=5]
设$x$，$y$满足约束条件$\begin{cases} x\ge 0 \\ y\ge 0 \\ x+y\le 3 \end{cases}$，则$z=x+y$的最大值为（   ）
\begin{choices}
  \item 0
  \item 1
  \item 2
  \item 3
\end{choices}
\begin{solution}
\textbf{答案：}D

\textbf{分析：}画出约束条件的可行域，利用目标函数的最优解求解目标函数的最大值即可。

\textbf{详解：}解：$x$，$y$满足约束条件$\begin{cases} x\ge 0 \\ y\ge 0 \\ x+y\le 3 \end{cases}$的可行域如图：$z=x+y$经过可行域的$A$时，目标函数取得最大值，由$\begin{cases} x=3 \\ y=0 \end{cases}$解得$A(3,0)$，所以$z=x+y$的最大值为：3。故选：D。
\end{solution}
\end{question}

\begin{question}[points=5]
函数$y=\frac{\sin x}{1+\cos x}$的部分图象大致为（        ）

\begin{center}
  \includegraphics[width=.38\linewidth]{question-8-1.jpg}
\end{center}
\begin{choices}
  \item 
  \item 
  \item 
  \item 
\end{choices}
\begin{solution}
\textbf{答案：}C

\textbf{分析：}判断函数的奇偶性排除选项，利用特殊值判断即可。

\textbf{详解：}解：函数$y=\frac{\sin x}{1+\cos x}$，可知函数是奇函数，排除选项B，当$x=\frac{\pi}{2}$时，$f(\frac{\pi}{2})=\frac{1}{1+0}=1$，排除A，$x=\pi$时，$f(\pi)=0$，排除D。故选：C。
\end{solution}
\end{question}

\begin{question}[points=5]
已知函数$f(x)=\ln x+\ln(2-x)$，则（           ）
\begin{choices}
  \item $f(x)$在$(0,2)$单调递增
  \item $f(x)$在$(0,2)$单调递减
  \item $y=f(x)$的图象关于直线$x=1$对称
  \item $y=f(x)$的图象关于点$(1,0)$对称
\end{choices}
\begin{solution}
\textbf{答案：}C

\textbf{分析：}由已知中函数$f(x)=\ln x+\ln(2-x)$，可得$f(x)=f(2-x)$，进而可得函数图象的对称性。

\textbf{详解：}解：$\because$函数$f(x)=\ln x+\ln(2-x)$，$\therefore f(2-x)=\ln(2-x)+\ln x$，即$f(x)=f(2-x)$，即$y=f(x)$的图象关于直线$x=1$对称，故选：C。
\end{solution}
\end{question}

\begin{question}[points=5]
如图程序框图是为了求出满足$3^n-2^n>1000$的最小偶数$n$，那么在$\boxed{}$和$\boxed{}$两个空白框中，可以分别填入（             ）

\begin{center}
  \includegraphics[width=.38\linewidth]{question-10-1.jpg}
\end{center}
\begin{choices}
  \item $A>1000$和$n=n+1$
  \item $A>1000$和$n=n+2$
  \item $A\le 1000$和$n=n+1$
  \item $A\le 1000$和$n=n+2$
\end{choices}
\begin{solution}
\textbf{答案：}D

\textbf{分析：}通过要求$A>1000$时输出且框图中在“否”时输出确定“$\boxed{}$”内不能输入“$A>1000$”，进而通过偶数的特征确定$n=n+2$。

\textbf{详解：}解：因为要求$A>1000$时输出，且框图中在“否”时输出，所以“$\boxed{}$”内不能输入“$A>1000$”，又要求$n$为偶数，且$n$的初始值为0，所以“$\boxed{}$”中$n$依次加2可保证其为偶数，所以D选项满足要求，故选：D。
\end{solution}
\end{question}

\begin{question}[points=5]
$\triangle ABC$的内角$A$，$B$，$C$的对边分别为$a$，$b$，$c$，已知$\sin B+\sin A(\sin C-\cos C)=0$，$a=2$，$c=\sqrt{2}$，则$C=$（    ）
\begin{choices}
  \item $\frac{\pi}{12}$
  \item $\frac{\pi}{6}$
  \item $\frac{\pi}{4}$
  \item $\frac{\pi}{3}$
\end{choices}
\begin{solution}
\textbf{答案：}B

\textbf{分析：}根据诱导公式和两角和的正弦公式以及正弦定理计算即可。

\textbf{详解：}解：$\sin B=\sin(A+C)=\sin A\cos C+\cos A\sin C$，$\because\sin B+\sin A(\sin C-\cos C)=0$，$\therefore\sin A\cos C+\cos A\sin C+\sin A\sin C-\sin A\cos C=0$，$\therefore\cos A\sin C+\sin A\sin C=0$，$\because\sin C\neq 0$，$\therefore\cos A=-\sin A$，$\therefore\tan A=-1$，$\because\frac{\pi}{2}<A<\pi$，$\therefore A=\frac{3\pi}{4}$，由正弦定理可得$\frac{a}{\sin A}=\frac{c}{\sin C}$，$\therefore\sin C=\frac{c\sin A}{a}$，$\because a=2$，$c=\sqrt{2}$，$\therefore\sin C=\frac{\sqrt{2}\times\frac{\sqrt{2}}{2}}{2}=\frac{1}{2}$，$\because a>c$，$\therefore C=\frac{\pi}{6}$，故选：B。
\end{solution}
\end{question}

\begin{question}[points=5]
设$A$，$B$是椭圆$C:\frac{x^2}{3}+\frac{y^2}{m}=1$长轴的两个端点，若$C$上存在点$M$满足$\angle AMB=120^\circ$，则$m$的取值范围是（            ）
\begin{choices}
  \item $(0,1]\cup[9,+\infty)$
  \item $(0,\frac{1}{3}]\cup[9,+\infty)$
  \item $(0,1]\cup[4,+\infty)$
  \item $(0,\frac{1}{3}]\cup[4,+\infty)$
\end{choices}
\begin{solution}
\textbf{答案：}A

\textbf{分析：}分类讨论，由要使椭圆C上存在点M满足$\angle AMB=120^\circ$，$\angle AMB\ge 120^\circ$，$\angle AMO\ge 60^\circ$，当假设椭圆的焦点在x轴上，$\tan\angle AMO=\frac{a}{b}\ge\tan60^\circ$，当即可求得椭圆的焦点在y轴上时，$m>3$，$\tan\angle AMO=\frac{b}{a}\ge\tan60^\circ=\sqrt{3}$，即可求得m的取值范围。

\textbf{详解：}解：假设椭圆的焦点在x轴上，则$0<m<3$时，设椭圆的方程为$\frac{x^2}{a^2}+\frac{y^2}{b^2}=1$（$a>b>0$），设$A(-a,0)$，$B(a,0)$，$M(x,y)$，$y>0$，则$a^2-x^2=\frac{a^2}{b^2}y^2$，$\angle MAB=\alpha$，$\angle MBA=\beta$，$\angle AMB=\gamma$，$\tan\alpha=\frac{y}{x+a}$，$\tan\beta=\frac{y}{a-x}$，则$\tan\gamma=\tan[\pi-(\alpha+\beta)]=-\tan(\alpha+\beta)=-\frac{\tan\alpha+\tan\beta}{1-\tan\alpha\tan\beta}=-\frac{\frac{y}{x+a}+\frac{y}{a-x}}{1-\frac{y^2}{a^2-x^2}}=-\frac{2ay}{a^2-x^2-y^2}=-\frac{2a}{b^2}\cdot\frac{y(\frac{a^2}{b^2}y^2+y^2)}{?}$（具体推导略），$\therefore\tan\gamma=-\frac{2ab^2}{a^2-b^2}\cdot\frac{1}{y}$，当y最大时，即$y=b$时，$\angle AMB$取最大值，$\therefore M$位于短轴的端点时，$\angle AMB$取最大值，要使椭圆C上存在点M满足$\angle AMB=120^\circ$，$\angle AMB\ge 120^\circ$，$\angle AMO\ge 60^\circ$，$\tan\angle AMO=\frac{a}{b}\ge\tan60^\circ=\sqrt{3}$，解得：$0<m\le 1$；当椭圆的焦点在y轴上时，$m>3$，当M位于短轴的端点时，$\angle AMB$取最大值，要使椭圆C上存在点M满足$\angle AMB=120^\circ$，$\angle AMB\ge 120^\circ$，$\angle AMO\ge 60^\circ$，$\tan\angle AMO=\frac{b}{a}\ge\tan60^\circ=\sqrt{3}$，解得：$m\ge 9$，$\therefore m$的取值范围是$(0,1]\cup[9,+\infty)$，故选A。
\end{solution}
\end{question}

\section*{二、填空题}
本题共4小题，每小题5分，共20分。

\begin{question}[points=5]
已知向量$\vec{a}=(-1,2)$，$\vec{b}=(m,1)$，若向量$\vec{a}+\vec{b}$与$\vec{a}$垂直，则$m=$ ．
\fillin[{$7$}]
\begin{solution}
\textbf{答案：}$7$

\textbf{分析：}利用平面向量坐标运算法则先求出$\vec{a}+\vec{b}$，再由向量$\vec{a}+\vec{b}$与$\vec{a}$垂直，利用向量垂直的条件能求出$m$的值。

\textbf{详解：}解：$\because$向量$\vec{a}=(-1,2)$，$\vec{b}=(m,1)$，$\therefore\vec{a}+\vec{b}=(-1+m,3)$，$\because$向量$\vec{a}+\vec{b}$与$\vec{a}$垂直，$\therefore(\vec{a}+\vec{b})\cdot\vec{a}=(-1+m)\times(-1)+3\times2=0$，解得$m=7$。故答案为：7。
\end{solution}
\end{question}

\begin{question}[points=5]
曲线$y=x^2+\frac{1}{x}$在点$(1,2)$处的切线方程为．
\fillin[{$x-y+1=0$}]
\begin{solution}
\textbf{答案：}$x-y+1=0$

\textbf{分析：}求出函数的导数，求出切线的斜率，利用点斜式求解切线方程即可。

\textbf{详解：}解：曲线$y=x^2+\frac{1}{x}$，可得$y'=2x-\frac{1}{x^2}$，切线的斜率为：$k=2-1=1$。切线方程为：$y-2=x-1$，即：$x-y+1=0$。故答案为：$x-y+1=0$。
\end{solution}
\end{question}

\begin{question}[points=5]
已知$\alpha\in(0,\frac{\pi}{2})$，$\tan\alpha=2$，则$\cos(\alpha-\frac{\pi}{4})=$ ．
\fillin[{$\frac{3\sqrt{10}}{10}$}]
\begin{solution}
\textbf{答案：}$\frac{3\sqrt{10}}{10}$

\textbf{分析：}根据同角的三角函数的关系求出$\sin\alpha=\frac{2\sqrt{5}}{5}$，$\cos\alpha=\frac{\sqrt{5}}{5}$，再根据两角差的余弦公式即可求出。

\textbf{详解：}解：$\because\alpha\in(0,\frac{\pi}{2})$，$\tan\alpha=2$，$\therefore\sin\alpha=2\cos\alpha$，$\because\sin^2\alpha+\cos^2\alpha=1$，解得$\sin\alpha=\frac{2\sqrt{5}}{5}$，$\cos\alpha=\frac{\sqrt{5}}{5}$，$\therefore\cos(\alpha-\frac{\pi}{4})=\cos\alpha\cos\frac{\pi}{4}+\sin\alpha\sin\frac{\pi}{4}=\frac{\sqrt{5}}{5}\times\frac{\sqrt{2}}{2}+\frac{2\sqrt{5}}{5}\times\frac{\sqrt{2}}{2}=\frac{3\sqrt{10}}{10}$，故答案为：$\frac{3\sqrt{10}}{10}$。
\end{solution}
\end{question}

\begin{question}[points=5]
已知三棱锥$S-ABC$的所有顶点都在球$O$的球面上，$SC$是球$O$的直径．若平面$SCA\perp$平面$SCB$，$SA=AC$，$SB=BC$，三棱锥$S-ABC$的体积为$9$，则球$O$的表面积为．
\fillin[{$36\pi$}]
\begin{solution}
\textbf{答案：}$36\pi$

\textbf{分析：}判断三棱锥的形状，利用几何体的体积，求解球的半径，然后求解球的表面积。

\textbf{详解：}解：三棱锥$S-ABC$的所有顶点都在球$O$的球面上，$SC$是球$O$的直径，若平面$SCA\perp$平面$SCB$，$SA=AC$，$SB=BC$，三棱锥$S-ABC$的体积为$9$，可知三角形$SBC$与三角形$SAC$都是等腰直角三角形，设球的半径为$r$，可得$\frac{1}{3}\times\frac{1}{2}\times2r\times r\times r=9$，解得$r=3$。球$O$的表面积为：$4\pi r^2=36\pi$。故答案为：$36\pi$。
\end{solution}
\end{question}

\section*{三、解答题}
共70分。解答应写出文字说明、证明过程或演算过程。第17～21题为必考题，每个试题考生都必须作答。第22、23题为选考题，考生根据要求作答。

\begin{question}[points=12]
记$S_n$为等比数列$\{a_n\}$的前$n$项和．已知$S_2=2$，$S_3=-6$．

（1）求$\{a_n\}$的通项公式；

（2）求$S_n$，并判断$S_{n+1}$，$S_n$，$S_{n+2}$是否成等差数列．
\begin{solution}
\textbf{答案：}（1）$a_n=(-2)^n$；（2）$S_n=-\frac{2}{3}[2+(-2)^{n+1}]$，$S_{n+1}+S_{n+2}=2S_n$，即成等差数列。

\textbf{分析：}（1）由题意可知$a_3=S_3-S_2=-8$，进而得到$a_1$，$a_2$，由$a_1+a_2=2$，列方程即可求得$q$及$a_1$，根据等比数列通项公式，即可求得$\{a_n\}$的通项公式；
（2）由（1）可知．利用等比数列前$n$项和公式，即可求得$S_n$，分别求得$S_{n+1}$，$S_{n+2}$，显然$S_{n+1}+S_{n+2}=2S_n$，则$S_{n+1}$，$S_n$，$S_{n+2}$成等差数列．

\textbf{详解：}解：（1）设等比数列$\{a_n\}$首项为$a_1$，公比为$q$，则$a_3=S_3-S_2=-6-2=-8$，则$a_1=\frac{a_3}{q^2}=\frac{-8}{q^2}$，$a_2=\frac{a_3}{q}=\frac{-8}{q}$，由$a_1+a_2=2$，$\frac{-8}{q^2}+\frac{-8}{q}=2$，整理得：$q^2+4q+4=0$，解得：$q=-2$，则$a_1=-2$，$a_n=(-2)(-2)^{n-1}=(-2)^n$，$\therefore\{a_n\}$的通项公式$a_n=(-2)^n$；
（2）由（1）可知：$S_n=\frac{a_1(1-q^n)}{1-q}=\frac{-2[1-(-2)^n]}{1-(-2)}=-\frac{2}{3}[1-(-2)^n]=-\frac{2}{3}[2+(-2)^{n+1}]$，则$S_{n+1}=-\frac{2}{3}[2+(-2)^{n+2}]$，$S_{n+2}=-\frac{2}{3}[2+(-2)^{n+3}]$，由$S_{n+1}+S_{n+2}=-\frac{2}{3}[2+(-2)^{n+2}]-\frac{2}{3}[2+(-2)^{n+3}]=-\frac{2}{3}[4+(-2)^{n+2}+(-2)^{n+3}]=-\frac{2}{3}[4+(-2)^{n+2}(1-2)]=-\frac{2}{3}[4-(-2)^{n+2}]=-\frac{2}{3}[4-4\times(-2)^n]=-\frac{2}{3}\times4[1-(-2)^n]=2\times(-\frac{2}{3}[1-(-2)^n])=2S_n$，即$S_{n+1}+S_{n+2}=2S_n$，$\therefore S_{n+1}$，$S_n$，$S_{n+2}$成等差数列．
\end{solution}
\end{question}

\begin{question}[points=12]
如图，在四棱锥$P-ABCD$中，$AB\parallel CD$，且$\angle BAP=\angle CDP=90^\circ$．

（1）证明：平面$PAB\perp$平面$PAD$；

（2）若$PA=PD=AB=DC$，$\angle APD=90^\circ$，且四棱锥$P-ABCD$的体积为$\frac{8}{3}$，求该四棱锥的侧面积．

\begin{center}
  \includegraphics[width=.38\linewidth]{question-18-1.jpg}
\end{center}
\begin{solution}
\textbf{答案：}（1）证明见解析；（2）$6+2\sqrt{3}$。

\textbf{分析：}（1）推导出$AB\perp PA$，$CD\perp PD$，从而$AB\perp PD$，进而$AB\perp$平面$PAD$，由此能证明平面$PAB\perp$平面$PAD$．（2）设$PA=PD=AB=DC=a$，取$AD$中点$O$，连结$PO$，则$PO\perp$底面$ABCD$，且$AD=\sqrt{2}a$，$PO=\frac{\sqrt{2}}{2}a$，由四棱锥$P-ABCD$的体积为$\frac{8}{3}$，求出$a=2$，由此能求出该四棱锥的侧面积．

\textbf{详解：}证明：（1）$\because$在四棱锥$P-ABCD$中，$\angle BAP=\angle CDP=90^\circ$，$\therefore AB\perp PA$，$CD\perp PD$，又$AB\parallel CD$，$\therefore AB\perp PD$，$\because PA\cap PD=P$，$\therefore AB\perp$平面$PAD$，$\because AB\subset$平面$PAB$，$\therefore$平面$PAB\perp$平面$PAD$．解：（2）设$PA=PD=AB=DC=a$，取$AD$中点$O$，连结$PO$，$\because PA=PD=AB=DC$，$\angle APD=90^\circ$，平面$PAB\perp$平面$PAD$，$\therefore PO\perp$底面$ABCD$，且$AD=\sqrt{a^2+a^2}=\sqrt{2}a$，$PO=\frac{\sqrt{2}}{2}a$，$\because$四棱锥$P-ABCD$的体积为$\frac{8}{3}$，由$AB\perp$平面$PAD$，得$AB\perp AD$，$\therefore V_{P-ABCD}=\frac{1}{3}\times \frac{1}{2}(AB+CD)\times AD\times PO=\frac{1}{3}\times \frac{1}{2}(a+a)\times\sqrt{2}a\times\frac{\sqrt{2}}{2}a=\frac{1}{3}a^3=\frac{8}{3}$，解得$a=2$，$\therefore PA=PD=AB=DC=2$，$AD=BC=2\sqrt{2}$，$PO=\frac{\sqrt{2}}{2}\times2=\sqrt{2}$，$\therefore PB=PC=\sqrt{PA^2+AB^2}=\sqrt{4+4}=2\sqrt{2}$，$\therefore$该四棱锥的侧面积：$S_{\text{侧}}=S_{\triangle PAD}+S_{\triangle PAB}+S_{\triangle PDC}+S_{\triangle PBC}=\frac{1}{2}\times2\times2+\frac{1}{2}\times2\times2+\frac{1}{2}\times2\times2+\frac{1}{2}\times2\sqrt{2}\times\sqrt{(2\sqrt{2})^2-(\frac{2\sqrt{2}}{2})^2}=2+2+2+\frac{1}{2}\times2\sqrt{2}\times\sqrt{8-2}=6+\frac{1}{2}\times2\sqrt{2}\times\sqrt{6}=6+\sqrt{12}=6+2\sqrt{3}$．
\end{solution}
\end{question}

\begin{question}[points=12]
为了监控某种零件的一条生产线的生产过程，检验员每隔$30\text{min}$从该生产线上随机抽取一个零件，并测量其尺寸（单位：$\text{cm}$）．下面是检验员在一天内依次抽取的$16$个零件的尺寸：



抽取次序 1 2 3 4 5 6 7 8

零件尺寸 9.95 10.12 9.96 9.96 10.01 9.92 9.98 10.04

抽取次序 9 10 11 12 13 14 15 16

零件尺寸 10.26 9.91 10.13 10.02 9.22 10.04 10.05 9.95



经计算得$\bar{x}=\frac{1}{16}\sum_{i=1}^{16}x_i=9.97$，$s=\sqrt{\frac{1}{16}\sum_{i=1}^{16}(x_i-\bar{x})^2}=\sqrt{\frac{1}{16}(\sum_{i=1}^{16}x_i^2-16\bar{x}^2)}\approx0.212$，$\sqrt{\sum_{i=1}^{16}(i-8.5)^2}\approx18.439$，$\sum_{i=1}^{16}(x_i-\bar{x})(i-8.5)=-2.78$，其中$x_i$为抽取的第$i$个零件的尺寸，$i=1,2,\ldots,16$．

（1）求$(x_i,i)$（$i=1,2,\ldots,16$）的相关系数$r$，并回答是否可以认为这一天生产的零件尺寸不随生产过程的进行而系统地变大或变小（若$|r|<0.25$，则可以认为零件的尺寸不随生产过程的进行而系统地变大或变小）．

（2）一天内抽检零件中，如果出现了尺寸在$(\bar{x}-3s,\bar{x}+3s)$之外的零件，就认为这条生产线在这一天的生产过程可能出现了异常情况，需对当天的生产过程进行检查．

（i）从这一天抽检的结果看，是否需对当天的生产过程进行检查？

（ii）在$(\bar{x}-3s,\bar{x}+3s)$之外的数据称为离群值，试剔除离群值，估计这条生产线当天生产的零件尺寸的均值与标准差．（精确到$0.01$）



附：样本$(x_i,y_i)$（$i=1,2,\ldots,n$）的相关系数$r=\frac{\sum_{i=1}^{n}(x_i-\bar{x})(y_i-\bar{y})}{\sqrt{\sum_{i=1}^{n}(x_i-\bar{x})^2}\sqrt{\sum_{i=1}^{n}(y_i-\bar{y})^2}}$，$\sqrt{0.008}\approx0.09$．

\begin{center}
  \includegraphics[width=.38\linewidth]{question-19-1.jpg}
\end{center}
\begin{solution}
\textbf{答案：}（1）$r\approx-0.18$，可以认为；（2）（i）需要对当天的生产过程进行检查；（ii）均值约为$10.02$，标准差约为$0.09$。

\textbf{分析：}（1）代入数据计算，比较$|r|$与$0.25$的大小作出结论；（2）（i）计算合格零件尺寸范围，得出结论；（ii）代入公式计算即可。

\textbf{详解：}解：（1）$r=\frac{\sum_{i=1}^{16}(x_i-\bar{x})(i-8.5)}{\sqrt{\sum_{i=1}^{16}(x_i-\bar{x})^2}\sqrt{\sum_{i=1}^{16}(i-8.5)^2}}=\frac{-2.78}{16\times0.212\times18.439}\approx-0.18$．$\because|r|<0.25$，$\therefore$可以认为这一天生产的零件尺寸不随生产过程的进行而系统地变大或变小．
（2）（i）$\bar{x}=9.97$，$s=0.212$，$\therefore$合格零件尺寸范围是$(9.334,10.606)$，显然第$13$号零件尺寸$9.22$不在此范围之内，$\therefore$需要对当天的生产过程进行检查．
（ii）剔除离群值后，剩下的数据平均值为$\frac{16\times9.97-9.22}{15}=10.02$，$\sum_{i=1}^{16}x_i^2=16\times0.212^2+16\times9.97^2=1591.134$，$\therefore$剔除离群值后样本方差为$\frac{1}{15}(1591.134-9.22^2-15\times10.02^2)=0.008$，$\therefore$剔除离群值后样本标准差为$\sqrt{0.008}\approx0.09$．
\end{solution}
\end{question}

\begin{question}[points=12]
设$A$，$B$为曲线$C:y=\frac{x^2}{4}$上两点，$A$与$B$的横坐标之和为$4$．

（1）求直线$AB$的斜率；

（2）设$M$为曲线$C$上一点，$C$在$M$处的切线与直线$AB$平行，且$AM\perp BM$，求直线$AB$的方程．
\begin{solution}
\textbf{答案：}（1）$1$；（2）$y=x+7$。

\textbf{分析：}（1）设$A(x_1,\frac{x_1^2}{4})$，$B(x_2,\frac{x_2^2}{4})$，运用直线的斜率公式，结合条件，即可得到所求；（2）设$M(m,\frac{m^2}{4})$，求出$y=\frac{x^2}{4}$的导数，可得切线的斜率，由两直线平行的条件：斜率相等，可得$m$，即有$M$的坐标，再由两直线垂直的条件：斜率之积为$-1$，可得$x_1,x_2$的关系式，再由直线$AB:y=x+t$与$y=\frac{x^2}{4}$联立，运用韦达定理，即可得到$t$的方程，解得$t$的值，即可得到所求直线方程．

\textbf{详解：}解：（1）设$A(x_1,\frac{x_1^2}{4})$，$B(x_2,\frac{x_2^2}{4})$为曲线$C:y=\frac{x^2}{4}$上两点，则直线$AB$的斜率为$k=\frac{\frac{x_2^2}{4}-\frac{x_1^2}{4}}{x_2-x_1}=\frac{1}{4}(x_1+x_2)=\frac{1}{4}\times4=1$；
（2）设直线$AB$的方程为$y=x+t$，代入曲线$C:y=\frac{x^2}{4}$，可得$x^2-4x-4t=0$，即有$x_1+x_2=4$，$x_1x_2=-4t$，再由$y=\frac{x^2}{4}$的导数为$y'=\frac{x}{2}$，设$M(m,\frac{m^2}{4})$，可得$M$处切线的斜率为$\frac{m}{2}$，由$C$在$M$处的切线与直线$AB$平行，可得$\frac{m}{2}=1$，解得$m=2$，即$M(2,1)$，由$AM\perp BM$可得，$k_{AM}\cdot k_{BM}=-1$，即为$\frac{\frac{1}{4}x_1^2-1}{x_1-2}\cdot\frac{\frac{1}{4}x_2^2-1}{x_2-2}=-1$，化为$x_1x_2+2(x_1+x_2)+20=0$，即为$-4t+8+20=0$，解得$t=7$．则直线$AB$的方程为$y=x+7$．
\end{solution}
\end{question}

\begin{question}[points=12]
已知函数$f(x)=e^x(e^x-a)-a^2x$．

（1）讨论$f(x)$的单调性；

（2）若$f(x)\ge0$，求$a$的取值范围．
\begin{solution}
\textbf{答案：}（1）当$a=0$时，$f(x)$在$\mathbb{R}$上单调递增；当$a>0$时，$f(x)$在$(-\infty,\ln a)$上单调递减，在$(\ln a,+\infty)$上单调递增；当$a<0$时，$f(x)$在$(-\infty,\ln(-\frac{a}{2}))$上单调递减，在$(\ln(-\frac{a}{2}),+\infty)$上单调递增；（2）$[-2\sqrt{e},1]$。

\textbf{分析：}（1）先求导，再分类讨论，根据导数和函数的单调性即可判断；（2）根据（1）的结论，分别求出函数的最小值，即可求出$a$的范围。

\textbf{详解：}解：（1）$f(x)=e^x(e^x-a)-a^2x=e^{2x}-ae^x-a^2x$，$\therefore f'(x)=2e^{2x}-ae^x-a^2=(2e^x+a)(e^x-a)$，①当$a=0$时，$f'(x)>0$恒成立，$\therefore f(x)$在$\mathbb{R}$上单调递增，②当$a>0$时，$2e^x+a>0$，令$f'(x)=0$，解得$x=\ln a$，当$x<\ln a$时，$f'(x)<0$，函数$f(x)$单调递减，当$x>\ln a$时，$f'(x)>0$，函数$f(x)$单调递增，③当$a<0$时，$e^x-a>0$，令$f'(x)=0$，解得$x=\ln(-\frac{a}{2})$，当$x<\ln(-\frac{a}{2})$时，$f'(x)<0$，函数$f(x)$单调递减，当$x>\ln(-\frac{a}{2})$时，$f'(x)>0$，函数$f(x)$单调递增，综上所述，当$a=0$时，$f(x)$在$\mathbb{R}$上单调递增，当$a>0$时，$f(x)$在$(-\infty,\ln a)$上单调递减，在$(\ln a,+\infty)$上单调递增，当$a<0$时，$f(x)$在$(-\infty,\ln(-\frac{a}{2}))$上单调递减，在$(\ln(-\frac{a}{2}),+\infty)$上单调递增，
（2）①当$a=0$时，$f(x)=e^{2x}>0$恒成立，②当$a>0$时，由（1）可得$f(x)_{\min}=f(\ln a)=-a^2\ln a\ge0$，$\therefore\ln a\le0$，$\therefore0<a\le1$，③当$a<0$时，由（1）可得：$f(x)_{\min}=f(\ln(-\frac{a}{2}))=\frac{a^2}{4}-a^2\ln(-\frac{a}{2})\ge0$，$\therefore\ln(-\frac{a}{2})\le\frac{1}{4}$，$\therefore -2\sqrt{e}\le a<0$，综上所述$a$的取值范围为$[-2\sqrt{e},1]$．
\end{solution}
\end{question}

\begin{question}[points=10]
在直角坐标系$xOy$中，曲线$C$的参数方程为$\begin{cases} x=3\cos\theta \\ y=\sin\theta \end{cases}$（$\theta$为参数），直线$l$的参数方程为$\begin{cases} x=a+4t \\ y=1-t \end{cases}$（$t$为参数）．

（1）若$a=-1$，求$C$与$l$的交点坐标；

（2）若$C$上的点到$l$距离的最大值为$\sqrt{17}$，求$a$．
\begin{solution}
\textbf{答案：}（1）$(3,0)$和$(-\frac{21}{25},\frac{24}{25})$；（2）$a=8$或$a=-16$。

\textbf{分析：}（1）将曲线C的参数方程化为标准方程，直线l的参数方程化为一般方程，联立两方程可以求得交点坐标；（2）曲线C上的点可以表示成$P(3\cos\theta,\sin\theta)$，$\theta\in[0,2\pi)$，运用点到直线距离公式可以表示出P到直线l的距离，再结合距离最大值为$\sqrt{17}$进行分析，可以求出a的值。

\textbf{详解：}解：（1）曲线C的参数方程为$\begin{cases} x=3\cos\theta \\ y=\sin\theta \end{cases}$（$\theta$为参数），化为标准方程是：$\frac{x^2}{9}+y^2=1$；$a=-1$时，直线l的参数方程化为一般方程是：$x+4y-3=0$；联立方程$\begin{cases} \frac{x^2}{9}+y^2=1 \\ x+4y-3=0 \end{cases}$，解得$\begin{cases} x=3 \\ y=0 \end{cases}$或$\begin{cases} x=-\frac{21}{25} \\ y=\frac{24}{25} \end{cases}$，所以椭圆C和直线l的交点为$(3,0)$和$(-\frac{21}{25},\frac{24}{25})$．
（2）l的参数方程$\begin{cases} x=a+4t \\ y=1-t \end{cases}$（t为参数）化为一般方程是：$x+4y-a-4=0$，椭圆C上的任一点P可以表示成$P(3\cos\theta,\sin\theta)$，$\theta\in[0,2\pi)$，所以点P到直线l的距离d为：$d=\frac{|3\cos\theta+4\sin\theta-a-4|}{\sqrt{1^2+4^2}}=\frac{|5\sin(\theta+\varphi)-a-4|}{\sqrt{17}}$，$\varphi$满足$\tan\varphi=\frac{3}{4}$，且d的最大值为$\sqrt{17}$．①当$-a-4\le0$时，即$a\ge-4$时，$|5\sin(\theta+\varphi)-a-4|\le|5-a-4|=|1-a|$，最大值为$\frac{|1-a|}{\sqrt{17}}$，由$\frac{|1-a|}{\sqrt{17}}=\sqrt{17}$得$|1-a|=17$，解得$a=18$或$a=-16$（舍去），$a=18$满足$a\ge-4$；②当$-a-4>0$时，即$a<-4$时，$|5\sin(\theta+\varphi)-a-4|\le|5-a-4|=|1-a|=1-a$（因为$a<-4$，$1-a>0$），最大值为$\frac{1-a}{\sqrt{17}}$，由$\frac{1-a}{\sqrt{17}}=\sqrt{17}$得$1-a=17$，解得$a=-16$，满足$a<-4$．综上$a=18$或$a=-16$．
\end{solution}
\end{question}

\begin{question}[points=10]
已知函数$f(x)=-x^2+ax+4$，$g(x)=|x+1|+|x-1|$．

（1）当$a=1$时，求不等式$f(x)\ge g(x)$的解集；

（2）若不等式$f(x)\ge g(x)$的解集包含$[-1,1]$，求$a$的取值范围．
\begin{solution}
\textbf{答案：}（1）$[-1,\frac{\sqrt{17}-1}{2}]$；（2）$[-1,1]$。

\textbf{分析：}（1）当$a=1$时，$f(x)=-x^2+x+4$，$g(x)=|x+1|+|x-1|=\begin{cases} 2x, & x>1 \\ 2, & -1\le x\le 1 \\ -2x, & x<-1 \end{cases}$，分$x>1$、$x\in[-1,1]$、$x\in(-\infty,-1)$三类讨论，结合$g(x)$与$f(x)$的单调性质即可求得$f(x)\ge g(x)$的解集为$[-1,\frac{\sqrt{17}-1}{2}]$；（2）依题意得：$-x^2+ax+4\ge2$在$[-1,1]$恒成立$\Leftrightarrow x^2-ax-2\le0$在$[-1,1]$恒成立，只需$\begin{cases} 1-a-2\le0 \\ 1+a-2\le0 \end{cases}$，解之即可得a的取值范围．

\textbf{详解：}解：（1）当$a=1$时，$f(x)=-x^2+x+4$，是开口向下，对称轴为$x=\frac{1}{2}$的二次函数，$g(x)=|x+1|+|x-1|=\begin{cases} 2x, & x>1 \\ 2, & -1\le x\le 1 \\ -2x, & x<-1 \end{cases}$，当$x\in(1,+\infty)$时，令$-x^2+x+4=2x$，解得$x=\frac{\sqrt{17}-1}{2}$，$g(x)$在$(1,+\infty)$上单调递增，$f(x)$在$(1,+\infty)$上单调递减，$\therefore$此时$f(x)\ge g(x)$的解集为$(1,\frac{\sqrt{17}-1}{2}]$；当$x\in[-1,1]$时，$g(x)=2$，$f(x)\ge f(-1)=2$．当$x\in(-\infty,-1)$时，$g(x)$单调递减，$f(x)$单调递增，且$g(-1)=f(-1)=2$．综上所述，$f(x)\ge g(x)$的解集为$[-1,\frac{\sqrt{17}-1}{2}]$；
（2）依题意得：$-x^2+ax+4\ge2$在$[-1,1]$恒成立，即$x^2-ax-2\le0$在$[-1,1]$恒成立，则只需$\begin{cases} (-1)^2-a\cdot(-1)-2\le0 \\ 1^2-a\cdot1-2\le0 \end{cases}$，解得$-1\le a\le1$，故a的取值范围是$[-1,1]$．
\end{solution}
\end{question}

